\section{Conclusion}
\label{sec:conclusion}
% Budget: 0.25 page.

Auto-labeling a honeypot's traffic and retraining a detector on it treats the label as free
because the source is a decoy. It is not free: the labelling policy is public, the attacker
chooses exactly what the decoy sees, and the honeypot becomes a write channel into the
defender's training set that the defender opened on purpose. We showed this channel is
exploitable on a modern, flow-based retraining loop with a single mechanism — a label that
contradicts its neighbourhood — that a fixed detection threshold reads as false positives and a
recalibrated one reads as lost detection, not two separate effects. The attack needs near-copy
mimicry fidelity and nothing else: no model access, no knowledge of the poison ratio, no cost
beyond producing traffic that already resembles what the network generates. A generic
label-consistency defense cannot see this poison because it agrees with the defender's own
policy; a one-line cap on the honeypot's share of the training set can, and beats every
alternative we compared it against without needing any of the information those alternatives
require. The honest version of this loop works — that was never in question — but the design
that makes it convenient, a public rule that trusts a decoy's traffic unconditionally, is the
same design that makes it attackable, and the fix costs almost nothing to add.
